
\section{Current ellis Protocol: Proof-Oriented Pseudocode}

\paragraph{Scope.}
This pseudocode records the current protocol described in prose.  Each operation
$o$ has data $(\mathit{type},\mathit{key},\mathit{value},\mathit{client},
\mathit{seqno})$.  A client may have several outstanding operations.  At the
invocation of $o$, $\preds{o}$ is a snapshot of all other operations then
outstanding at that client.  The relation $p\in\preds{o}$ is the protocol's
cross-shard predecessor relation.

\paragraph{Messages.}
The coordination messages are
$\textsc{CoordReq}(p,o,L_o)$,
$\textsc{Arrival}(p,o,\ats{p})$, and
$\textsc{FinalAck}(p,o)$.  The last message contains no timestamp; it asserts
only that $p$'s final timestamp has been fixed.  Messages are idempotent and
identified by operation IDs.  Duplicate messages and duplicate VR replies do
not change state.

\paragraph{Shard state.}
Each shard leader maintains a monotonically increasing $\mathit{ShardTS}$, an
unordered durable bag produced by the first VR round, one pending ordering
queue $\poq{k}$ for every key $k$, and one shard-wide ordered commit log
$\clog$.  A POQ is sorted by the operation's current $\fts{}$ value.  The
commit log is ordered by its append indices and is also the shard's execution
order.

\paragraph{Ordering facts exposed to a correctness proof.}
The algorithm makes three relations explicit.  First,
$p\rightarrow_{\mathrm{inv}}o$ is recorded by the client when $p$ precedes
$o$ in that client's invocation order; outstanding such operations form
$\preds{o}$.  Second, operations on one object execute in the order induced by
their positions in their shard's commit log.  Third, a response is sent only
after execution, so if $p$ returns before $o$ is invoked, then
$p\rightarrow_{\mathrm{rt}}o$.  The intended later proof must show that the
union of per-object execution order, client invocation order, and real-time
order is acyclic.

\paragraph{Unresolved semantic obligations.}
The red annotations are deliberately not filled in by guesses.  In addition,
the prose increments $\mathit{ShardTS}$ only when assigning an arrival
timestamp.  It does not advance the counter when a much larger final timestamp
is chosen.  Consequently, a later same-key operation can receive an arrival
timestamp below an earlier operation's already-fixed final timestamp.  The
protocol or proof must explain why this cannot invalidate the POQ decision, or
must add a clock-advance rule.  The timing and durability of successor
registration across leader changes also remain unspecified.

